% Part: many-valued-logic
% Chapter: syntax-and-semantics
% Section: connectives

\documentclass[../../../include/open-logic-section]{subfiles}

\begin{document}

\olfileid[hi]{mvl}{syn}{con}

\olsection{भाषाएँ और संयोजक}

चिरसम्मत प्रतिज्ञप्ति तर्कशास्त्र और अनेक अन्य तर्कशास्त्र
\emph{प्रतिज्ञप्तिपरक अचरों} तथा \emph{संयोजकों} के नियत भंडार का उपयोग
करते हैं। उदाहरणार्थ, हम निम्न को आदिम मानते हैं:
\begin{enumerate}
\tagitem{prvFalse}{!!{falsity} के लिए प्रतिज्ञप्तिपरक अचर~$\lfalse$।}{}
\tagitem{prvTrue}{!!{truth} के लिए प्रतिज्ञप्तिपरक अचर~$\ltrue$।}{}
\item तार्किक संयोजक:
  \startycommalist
  \iftag{prvNot}{\ycomma $\lnot$ (निषेध)}{}%
  \iftag{prvAnd}{\ycomma $\land$ (संयोजन)}{}%
  \iftag{prvOr}{\ycomma $\lor$ (वियोजन)}{}%
  \iftag{prvIf}{\ycomma $\lif$ (!!{conditional})}{}%
  \iftag{prvIff}{\ycomma $\liff$ (!!{biconditional})}{}%
\end{enumerate}
\iftag{defNot,defOr,defAnd,defIf,defIff,defTrue,defFalse,defEx,defAll}{%
ऊपर के आदिम संयोजकों के अतिरिक्त हम संक्षिप्त रूपों के रूप में परिभाषित
चिह्न भी उपयोग करते हैं, जैसे
\startycommalist
  \iftag{defNot}{\ycomma $\lnot$ (निषेध)}{}%
  \iftag{defAnd}{\ycomma $\land$ (संयोजन)}{}%
  \iftag{defOr}{\ycomma $\lor$ (वियोजन)}{}%
  \iftag{defIf}{\ycomma $\lif$ (!!{conditional})}{}%
  \iftag{defIff}{\ycomma $\liff$ (!!{biconditional})}{}%
  \iftag{defFalse}{\ycomma $\lfalse$ (!!{falsity})}{}%
  \iftag{defTrue}{\ycomma $\ltrue$ (!!{truth})}.}{}

यही संयोजक बहुमूल्य तर्कशास्त्रों में भी उपयोग होते हैं। किंतु एक ही
तर्कशास्त्र में, मसलन, संयोजन के भिन्न रूप सम्मिलित करना प्रायः उपयोगी
होता है; उन्हें अलग रखने के लिए भिन्न चिह्न चाहिए होंगे। कुछ बहुमूल्य
तर्कशास्त्रों में ऐसे संयोजक भी होते हैं जिनका चिरसम्मत तर्कशास्त्र में कोई
तुल्य नहीं। अतः हम सामान्य से कुछ अधिक व्यापक होंगे।

\begin{defn}
  किसी \emph{प्रतिज्ञप्तिपरक भाषा} में \emph{संयोजकों} का समुच्चय $\Lang
  L$ होता है। प्रत्येक संयोजक $\star$ की कोई \emph{स्थानिकता} होती है;
  स्थानिकता~$n$ वाला संयोजक \emph{$n$-स्थानी} कहलाता है। स्थानिकता~$0$ के
  संयोजक \emph{अचर}, स्थानिकता~$1$ के संयोजक \emph{एक-स्थानी}, और
  स्थानिकता~$2$ के संयोजक \emph{दो-स्थानी} भी कहलाते हैं।
\end{defn}

\begin{ex}
  प्रतिज्ञप्ति तर्कशास्त्र की मानक भाषा $\Lang L_0$ में निम्न संयोजक (उनकी
  स्थानिकताओं सहित) हैं:
  $\lfalse$~($0$)
  $\lnot$~($1$),
  $\land$~($2$),
  $\lor$~($2$),
  $\lif$~($2$)। हमारे विचाराधीन अधिकांश तर्कशास्त्र इस भाषा का उपयोग
  करेंगे। कुछ तर्कशास्त्र परंपरा और रूढ़ि के कारण कुछ संयोजकों के लिए भिन्न
  चिह्न उपयोग करते हैं। उदाहरणार्थ, गुणनफल तर्कशास्त्र में संयोजन का चिह्न
  प्रायः~$\land$ के बदले $\odot$ होता है। कभी-कभी नया संक्रियक जोड़ना
  सुविधाजनक होता है, जैसे निर्धार्यता संक्रियक $\triangle$ ($1$-स्थानी)।
\end{ex}

\end{document}
